\section{Source work and its iterated time integrals}
\label{r:sec:source-primitives}

The source work can be controlled further before invoking any higher
velocity rectangle. For the base temporal row, write \(W_s=W_{0,s}\)
and put \(h_s=\Lambda^{2s}\Pp f\). Moving the spatial derivatives onto
the prescribed field gives \(W_s=\ip{u}{h_s}\). This pairing measures
how the fluid responds to a smooth spatial weight. Kinetic energy alone bounds this pairing; no higher spatial norm
of the velocity enters the estimate.

The energy balance implies
\[
 \norm{u(t)}_2\le K_T,
 \qquad K_T=\norm{u_0}_2+\int_0^T\norm{f(t)}_2\dd t.
\]
Let \(S_s=\operatorname{Sym}\nabla h_s\). Substitution of the
momentum equation into the derivative of \(\ip{u}{h_s}\) gives
\begin{equation}
 W_s'=\ip{u}{(\partial_t+\nu\Delta)h_s}
 +\int_{\R^3}u_i u_j(S_s)_{ij}\dd x+\ip{f}{h_s}.
 \label{r:eq:source-first}
\end{equation}
Incompressibility removes the pressure pairing; integration by parts
moves the transport and viscous derivatives onto \(h_s\).
Using the matrix operator norm for \(S_s\), define
\begin{equation}
 \begin{aligned}
 \mathsf B_s&:=K_T\norm{h_s}_{L^\infty_tL^2_x},\\
 \mathsf L_s&:=K_T\norm{(\partial_t+\nu\Delta)h_s}_{L^\infty_tL^2_x}
 +K_T^2\norm{S_s}_{L^\infty_{t,x}}\\
 &\quad+\norm{f}_{L^\infty_tL^2_x}\norm{h_s}_{L^\infty_tL^2_x}.
 \end{aligned}
 \label{r:eq:source-work-constants}
\end{equation}
Every fixed constant is finite under the source assumptions, and
\(\norm{W_s}_\infty\le\mathsf B_s\),
\(\norm{W_s'}_\infty\le\mathsf L_s\). The projection need not preserve
rapid spatial decrease; its boundedness on Sobolev spaces supplies the
norms used here.

These two bounds control an integrated form of the complete source
correction at every spatial height. Define
\begin{equation}
 \begin{aligned}
 \mathcal S_{s,k}&:=\sum_{j=0}^{k-1}(-2\nu)^j
 \partial_t^{k-1-j}W_{s+j},\qquad k\ge2,\\
 I^0&:=\mathrm{Id},\qquad
 (I^a g)(t):=\frac1{(a-1)!}\int_0^t(t-v)^{a-1}g(v)\dd v,
 \quad a\ge1.
 \end{aligned}
 \label{r:eq:source-primitive-def}
\end{equation}

\begin{proposition}[Integrated source correction]
\label{r:prop:source-primitive}
For each fixed \(s\ge0\) and integer \(k\ge2\),
\(I^{k-2}\mathcal S_{s,k}\) is bounded on \((0,T)\) by the datum and
source. If \(u_0=0\) and the force vanishes near \(t=0\), then
\begin{equation}
 \begin{aligned}
 I^{k-2}\mathcal S_{s,k}
 &=W_s'-2\nu W_{s+1}
 +\sum_{j=2}^{k-1}(-2\nu)^jI^{j-1}W_{s+j},\\
 \norm{I^{k-2}\mathcal S_{s,k}}_\infty
 &\le\mathsf L_s+
 \sum_{j=1}^{k-1}\frac{(2\nu)^jT^{j-1}}{(j-1)!}\mathsf B_{s+j}.
 \end{aligned}
 \label{r:eq:source-primitive-rest}
\end{equation}
No source amplitude restriction or higher velocity estimate is used.
\end{proposition}

\begin{proof}
Write \(D=\partial_t\) and \(m=k-2\). Repeated integration gives
\begin{equation}
 I^mD^dW=
 \begin{cases}
 D^{d-m}W-\displaystyle\sum_{\ell=0}^{m-1}
 W^{(d-m+\ell)}(0)\dfrac{t^\ell}{\ell!},&d\ge m,\\[5pt]
 I^{m-d}W-\displaystyle\sum_{\ell=0}^{d-1}
 W^{(\ell)}(0)\dfrac{t^{m-d+\ell}}{(m-d+\ell)!},&d<m.
 \end{cases}
 \label{r:eq:primitive-initial-terms}
\end{equation}
Empty sums vanish. In summand \(j\) of \(\mathcal S_{s,k}\), set
\(d=k-1-j=m+1-j\). The cases \(j=0,1\) leave \(W_s'\) and
\(-2\nu W_{s+1}\); the other cases leave the integrals in
\eqref{r:eq:source-primitive-rest}. Each initial correction is a polynomial
whose coefficients are determined by
\[
 W_s^{(\ell)}(0)=\sum_{b=0}^{\ell}\binom\ell b
 \ip{a_b}{\partial_t^{\ell-b}h_s(0)}.
\]
They are finite by \eqref{r:eq:initial}. Equation
\eqref{r:eq:primitive-initial-terms} specifies every correction for general
datum. The inequality \(\norm{I^aW}_\infty\le T^a\norm{W}_\infty/a!\)
proves the bound. For rest datum and forcing that vanishes initially,
local uniqueness gives a zero initial velocity history. All initial
corrections vanish, yielding \eqref{r:eq:source-primitive-rest}.
\end{proof}

The first instance bounds \(W_s'-2\nu W_{s+1}\) directly. The next
bounds the time integral of
\(W_s''-2\nu W_{s+1}'+4\nu^2W_{s+2}\). At higher orders the bound
uses exactly \(k-2\) integrations. It controls the signed source
contribution in that coordinate, rather than the absolute integral of
each differentiated term. This specifies which part of the completed
height calculation is bounded before using the velocity intersection.

\subsection{Source absorption in every temporal row}
\label{r:sec:all-row-source-absorption}

The kinetic estimate above bounds the pairing with \(u\). At a higher
temporal row, source work pairs the prescribed \(F_r\) with the response
\(V_r\). Viscous dissipation controls one additional spatial derivative
of \(V_r\). Pairing that derivative with one lower source derivative
allows part of the dissipation to absorb the source work.
Since \(V_r\) is divergence free,
\begin{equation}
 \begin{aligned}
 W_{r,s}
 &=\operatorname{Re}\ip{\Lambda^sV_r}{\Lambda^s\Pp F_r}\\
 &=\operatorname{Re}\ip{\Lambda^{s+1}V_r}
                         {\Lambda^{s-1}\Pp F_r}.
 \end{aligned}
 \label{r:eq:all-row-source-duality}
\end{equation}
For the prescribed rapidly decreasing force, the second factor is finite
at every \(s\ge0\). If \(s<1\), its Fourier weight has a singularity
at zero, but \(\widehat F_r\) is bounded there and the three-dimensional
radial integral is \(\int_0^1\rho^{2s}\dd\rho<\infty\).
The smoothness of the force controls the high frequencies. All these
bounds are uniform on the prescribed time interval.

Young's inequality now allocates one quarter of the viscous dissipation
to source work:
\begin{equation}
 \begin{aligned}
 |W_{r,s}|
 &\le\frac\nu4\norm{\Lambda^{s+1}V_r}_2^2
       +\frac1\nu\norm{\Lambda^{s-1}\Pp F_r}_2^2\\
 &=\frac\nu2E_{r,s+1}
       +\frac1\nu\norm{\Lambda^{s-1}\Pp F_r}_2^2.
 \end{aligned}
 \label{r:eq:all-row-source-absorption}
\end{equation}
Substitution into the complete height identity gives
\begin{equation}
 \partial_tE_{r,s}+\frac{3\nu}{2}E_{r,s+1}
 \le\mathcal P_{r,s}
       +\frac1\nu\norm{\Lambda^{s-1}\Pp F_r}_2^2.
 \label{r:eq:all-row-source-retained-production}
\end{equation}
The nonlinear production retains its sign and its full binomial sum.
The displayed source norm has a finite time integral determined by the
prescribed force. An absolute integral of \(W_{r,s}\) also requires
control of the dissipation appearing in
\eqref{r:eq:all-row-source-absorption}.

The argument can be written for the broader source class
\(\bigcap_{r,m}H^r_tH^m_x\), without a spatial-decay assumption.
Use \(J=(1-\Delta)^{1/2}\) and set
\[
 A_{r,s}=\tfrac12\norm{J^sV_r}_2^2,
 \qquad D_{r,s}=\norm{\nabla J^sV_r}_2^2,
 \qquad
 \mathcal P^J_{r,s}=-\operatorname{Re}\ip{J^sV_r}{J^sB_r}.
\]
Pressure again pairs to zero. The identity
\(\norm{J^{s+1}V_r}_2^2=D_{r,s}+2A_{r,s}\)
and Young's inequality yield
\begin{equation}
 A'_{r,s}+\frac\nu2D_{r,s}
 \le\mathcal P^J_{r,s}+\nu A_{r,s}
       +\frac1{2\nu}\norm{J^{s-1}\Pp F_r}_2^2.
 \label{r:eq:inhomogeneous-source-absorption}
\end{equation}
For \(s<1\), the multiplier defining \(J^{s-1}\) is bounded at
zero. Multiplying by \(e^{-\nu t}\) and integrating controls the
source contribution on every strict window with constants independent
of the time cutoff. The resulting inequality still contains the
signed integral of \(\mathcal P^J_{r,s}\).

\subsection{The source primitive at every temporal row}
\label{r:sec:all-row-primitives}

For \(h_s=\Lambda^{2s}\Pp f\), the kinetic bound also controls
the full direct source current at every finite temporal row.
Set

\[
M_{s,k}(t)=\operatorname{Re}\langle u(t),\partial_t^kh_s(t)\rangle.
\]

The exact binomial transfer is

\[
W_{r,s}
=\sum_{b=0}^r(-1)^b\binom rb
\partial_t^{r-b}M_{s,r+b}.
\]

To verify all indices and signs, expand the derivative on the right. The coefficient of
$\langle V_a,\partial_t^{2r-a}h_s\rangle$ is

\[
\sum_{b=0}^{r-a}(-1)^b\binom rb\binom{r-b}{a}
=\binom ra\sum_{b=0}^{r-a}(-1)^b\binom{r-a}{b}.
\]

It is zero for $a<r$ and one for $a=r$. The surviving term is
$\langle V_r,\partial_t^rh_s\rangle=W_{r,s}$.

The momentum calculation in \eqref{r:eq:source-first} also applies to
every prescribed test field \(g_{s,k}=\partial_t^kh_s\). It gives
\begin{equation}
\begin{aligned}
 M_{s,k}'={}&\operatorname{Re}\ip{u}{(\partial_t+\nu\Delta)g_{s,k}}\\
 &+\int_{\R^3}u_i u_j(\operatorname{Sym}\nabla g_{s,k})_{ij}\dd x
 +\operatorname{Re}\ip{f}{g_{s,k}}.
 \label{r:eq:all-row-test-derivative}
\end{aligned}
\end{equation}
Thus both the test pairing and its first time derivative are bounded
by kinetic energy and prescribed source norms:
\begin{equation}
\begin{aligned}
 \norm{M_{s,k}}_\infty&\le\mathsf B_{s,k}
 :=K_T\norm{g_{s,k}}_{L^\infty_tL^2_x},\\
 \norm{M_{s,k}'}_\infty&\le\mathsf L_{s,k}
 :=K_T\norm{(\partial_t+\nu\Delta)g_{s,k}}_{L^\infty_tL^2_x}\\
 &\quad+K_T^2\norm{\operatorname{Sym}\nabla g_{s,k}}_{L^\infty_{t,x},\mathrm{op}}
 +\norm{f}_{L^\infty_tL^2_x}\norm{g_{s,k}}_{L^\infty_tL^2_x}.
 \label{r:eq:all-row-test-bounds}
\end{aligned}
\end{equation}
The additional derivative bound removes two integrations from the
source primitive at every temporal row. Write the complete source
correction as
\[
 \mathcal F_{r,s,n}:=\sum_{j=0}^{n-1}(-2\nu)^j
 \partial_t^{n-1-j}W_{r,s+j},\qquad n\ge1.
\]

\begin{proposition}[Source correction with two fewer integrations]
\label{r:prop:all-row-sharp-primitive}
For every fixed \(r\ge0\), \(n\ge1\), and \(s\ge0\), the complete
source correction satisfies
\begin{equation}
 \norm{I^{\max(n+r-2,0)}\mathcal F_{r,s,n}}_{L^\infty(0,T)}<\infty.
 \label{r:eq:all-row-sharp-primitive}
\end{equation}
The bound uses \(\nu,T\), the kinetic bound, prescribed force norms,
and finitely many initial jets. No restriction on the force amplitude
or higher velocity norm is required.
\end{proposition}

\begin{proof}
Substitution of the binomial transfer into the complete source
correction gives
\begin{equation}
 \mathcal F_{r,s,n}
 =\sum_{j=0}^{n-1}\sum_{b=0}^r(-2\nu)^j(-1)^b\binom rb
 \partial_t^{n+r-1-j-b}M_{s+j,r+b}.
 \label{r:eq:all-row-source-transfer}
\end{equation}
When \((r,n)=(0,1)\), this is \(M_{s,0}\), already bounded by
\eqref{r:eq:all-row-test-bounds}. Otherwise put \(q=n+r-2\). For
\(0\le v\le q+1\), repeated integration gives
\begin{equation}
 I^q\partial_t^{q+1-v}M=\mathcal J_{q,v}[M],\qquad
 \mathcal J_{q,v}[M]=
 \begin{cases}
 M'-\displaystyle\sum_{\ell=0}^{q-1}
 M^{(\ell+1)}(0)\dfrac{t^\ell}{\ell!},&v=0,\\[7pt]
 M-\displaystyle\sum_{\ell=0}^{q-1}
 M^{(\ell)}(0)\dfrac{t^\ell}{\ell!},&v=1,\\[7pt]
 I^{v-1}M-\displaystyle\sum_{\ell=0}^{q-v}
 M^{(\ell)}(0)\dfrac{t^{v-1+\ell}}{(v-1+\ell)!},&v\ge2.
 \end{cases}
 \label{r:eq:all-row-sharp-integration}
\end{equation}
Empty sums vanish. In summand \((j,b)\) of
\eqref{r:eq:all-row-source-transfer}, take \(v=j+b\). The first two
branches are bounded by \(\mathsf L_{s+j,r+b}\) and
\(\mathsf B_{s+j,r+b}\), together with their initial polynomials.
For the third branch,
\[
 \norm{I^{v-1}M_{s,k}}_\infty
 \le\frac{T^{v-1}}{(v-1)!}\mathsf B_{s,k}.
\]
Every initial coefficient is finite: the product rule writes it as
\[
 M_{s,k}^{(\ell)}(0)
 =\sum_{c=0}^{\ell}\binom\ell c
 \operatorname{Re}\ip{V_c(0)}{g_{s,k+\ell-c}(0)},
\]
and \eqref{r:eq:initial} determines each \(V_c(0)\) from the smooth
datum and prescribed force jet. There are finitely many summands for
each fixed \((r,s,n)\), proving \eqref{r:eq:all-row-sharp-primitive}.
\end{proof}

For the nonnegative integrated height identity below, it remains useful
to write the source correction after \(n+r\) integrations explicitly.
Apply $I^{n+r}$ to \(\mathcal F_{r,s,n}\). For $0\le j<n$ and $0\le b\le r$, put $q=j+b+1$. The integration formula \eqref{r:eq:primitive-initial-terms}, with total integration order $n+r$, gives

\[
\begin{aligned}
&I^{n+r}
\partial_t^{n-1-j+r-b}M_{s+j,r+b}(t)\\
&\quad=I^qM_{s+j,r+b}(t)
-\sum_{\ell=0}^{n+r-2-j-b}
M_{s+j,r+b}^{(\ell)}(0)
\frac{t^{q+\ell}}{(q+\ell)!}.
\end{aligned}
\]

An upper index of \(-1\) means an empty sum. Multiplying the preceding identity by
$(-2\nu)^j(-1)^b\binom rb$ and summing gives the entire source current, with no terminal derivative.

\Needspace{10\baselineskip}
For an explicit bound, set

\[
S_{s,k}=\sup_{0\le t\le T}\|\partial_t^kh_s(t)\|_2.
\]

Then $|M_{s,k}|\le K_TS_{s,k}$ and
\[
\sup_{t<T}|I^{n+r}\mathcal F_{r,s,n}(t)|\le B_{r,s,n},
\]
where

\[
\begin{aligned}
B_{r,s,n}:=
\sum_{j=0}^{n-1}\sum_{b=0}^r
(2\nu)^j\binom rb
\Bigg[
&K_TS_{s+j,r+b}\frac{T^{j+b+1}}{(j+b+1)!}\\
&+\sum_{\ell=0}^{n+r-2-j-b}
|M_{s+j,r+b}^{(\ell)}(0)|
\frac{T^{j+b+1+\ell}}{(j+b+1+\ell)!}
\Bigg]<\infty .
\end{aligned}
\]

Every initial jet is finite by the forced recurrence and smooth initial data. This establishes finite integrated direct forcing in every temporal row and at every spatial height, separately for each finite $r,s,n$.

Write
\(\mathcal N_{r,s,n}=\sum_{j=0}^{n-1}(-2\nu)^j
\partial_t^{n-1-j}\mathcal P_{r,s+j}\).
For odd \(n\), integrating \eqref{r:eq:completed} $n+r$ times gives

\[
\begin{aligned}
I^rE_{r,s}
+(2\nu)^nI^{n+r}E_{r,s+n}
={}&I^rT_{r,s,n}
+I^{n+r}\mathcal N_{r,s,n}\\
&+I^{n+r}\mathcal F_{r,s,n},
\end{aligned}
\]

where
$T_{r,s,n}(t)=\sum_{\ell=0}^{n-1}E_{r,s}^{(\ell)}(0)t^\ell/\ell!$.
Both left-hand terms are nonnegative. Thus divergence of $I^rE_{r,s}$ requires an unbounded upper value of $I^{n+r}\mathcal N_{r,s,n}$ at every odd $n$. For $r>0$, this conclusion requires divergence of the indicated time integral; pointwise divergence of $E_{r,s}$ alone is insufficient.

\subsection{Source work in the first two temporal rows}
\label{r:app:recovered-source-work}
The source work has a useful bound before any higher velocity rectangle is
read off. This includes mixed spatial derivatives and the first temporal
row. Keep the actual forced velocity, set \(V_r=\partial_t^ru\) and
\(F_r=\partial_t^rf\), and write
\[
 U_*:=\|u_0\|_2+\int_0^T\|f(t)\|_2\dd t.
\]
The kinetic identity gives \(\sup_{t<T}\|u(t)\|_2\le U_*\).
For a spatial multi-index \(\alpha\) and a real height \(s\ge0\), introduce
the nonnegative self-adjoint multiplier
\[
 M_{\alpha,s}=(-1)^{|\alpha|}\partial_x^{2\alpha}\Lambda^{2s},
 \qquad
 \widehat{M_{\alpha,s}v}(\xi)
   =\xi^{2\alpha}|\xi|^{2s}\widehat v(\xi).
\]
Spatial self-adjointness gives the exact source pairing
\begin{equation}
 W_{r,\alpha,s}
 :=\ip{\Lambda^s\partial^\alpha V_r}
          {\Lambda^s\partial^\alpha F_r}
 =\ip{V_r}{M_{\alpha,s}F_r}.
 \label{r:eq:recovered-mixed-source-work}
\end{equation}
All pairings in this section mean their real parts. At \(r=0\), every
spatial derivative can thus be carried by the prescribed source:
\begin{equation}
 \int_0^T|W_{0,\alpha,s}|\dd t
 \le U_*\int_0^T\|M_{\alpha,s}f\|_2\dd t<\infty.
 \label{r:eq:recovered-source-row-zero}
\end{equation}
The constant is finite at each chosen order. This estimate consumes no
viscous dissipation and requires no higher velocity norm.

For \(r=1\), let \(h=M_{\alpha,s}\Pp F_1\). This is a prescribed smooth
solenoidal field. Pairing the original momentum equation with \(h\),
moving spatial derivatives to \(h\), and using incompressibility gives
\begin{equation}
 W_{1,\alpha,s}
 =\nu\ip{u}{\Delta h}
  +\int_{\mathbb R^3}u_i u_j(\operatorname{Sym}\nabla h)_{ij}\dd x
  +\ip{f}{h}.
 \label{r:eq:recovered-source-row-one-identity}
\end{equation}
Pressure contributes zero because \(\nabla\cdot h=0\). The transport term
has the displayed positive sign after integration by parts; contraction
with \(u\otimes u\) removes the antisymmetric part of \(\nabla h\).
Consequently,
\begin{align}
 \int_0^T|W_{1,\alpha,s}|\dd t
 \le{}&\nu U_*\int_0^T\|\Delta h\|_2\dd t\notag\\
 &+U_*^2\int_0^T
       \|\operatorname{Sym}\nabla h\|_{L^\infty_x,\mathrm{op}}\dd t
   +\int_0^T\|f\|_2\|h\|_2\dd t<\infty.
 \label{r:eq:recovered-source-row-one-bound}
\end{align}
Every quantity on the right is fixed by the datum, viscosity, and source.
Sobolev boundedness of \(\Pp\), followed by an adequately high spatial
embedding, controls these norms. No spatial support or rapid-decay
property of \(\Pp f\) is required. The estimate imposes no restriction on
force amplitude.

The higher temporal rows have an exact finite reduction. For \(r\ge1\),
put
\[
 B_{r,\alpha,s}
 =\sum_{j=0}^{r-1}(-1)^j
   \ip{V_{r-1-j}}{M_{\alpha,s}F_{r+j}}.
\]
Differentiating this sum cancels adjacent terms and leaves
\[
 B_{r,\alpha,s}'
 =W_{r,\alpha,s}
  +(-1)^{r-1}\ip{u}{M_{\alpha,s}F_{2r}}.
\]
Hence
\begin{equation}
 \int_0^tW_{r,\alpha,s}\dd v
 =B_{r,\alpha,s}(t)-B_{r,\alpha,s}(0)
  +(-1)^r\int_0^t\ip{u}{M_{\alpha,s}F_{2r}}\dd v.
 \label{r:eq:recovered-source-temporal-telescope}
\end{equation}
The terminal pairings involve temporal velocity rows strictly below
\(r\). If \(K_j=\sup_{t<T}\|V_j(t)\|_2\) for \(j<r\), the explicit bound
is
\begin{align}
 \sup_{t<T}\left|\int_0^tW_{r,\alpha,s}\dd v\right|
 \le{}&2\sum_{j=0}^{r-1}K_{r-1-j}
           \sup_{t\le T}\|M_{\alpha,s}F_{r+j}(t)\|_2\notag\\
 &+U_*\int_0^T\|M_{\alpha,s}F_{2r}\|_2\dd t.
 \label{r:eq:recovered-source-temporal-bound}
\end{align}
At \(r=1\), only the kinetic bound is needed. At larger \(r\), the stated
lower temporal norms enter explicitly; the datum-generated intersection
supplies them in the principal argument. Formula
\eqref{r:eq:recovered-source-temporal-bound} controls a signed integral,
while \eqref{r:eq:recovered-source-row-one-bound} controls its absolute
integral at the first temporal row.

For completeness, with
\(E_{r,\alpha,s}=\|\Lambda^s\partial^\alpha V_r\|_2^2/2\),
the corresponding energy identity can be written
\[
 \partial_t(E_{r,\alpha,s}-B_{r,\alpha,s})
 =-\nu\|\Lambda^{s+1}\partial^\alpha V_r\|_2^2
   +\mathcal P_{r,\alpha,s}
   +(-1)^r\ip{u}{M_{\alpha,s}F_{2r}}.
\]
Here \(\mathcal P_{r,\alpha,s}\) retains the complete differentiated
nonlinear row and its pressure completion. This identity changes the
recorded scalar by an explicitly bounded pairing when the lower rows are
bounded. It leaves the underlying velocity and its equation intact.

\subsection{Source insertions in the velocity rows}
Write
\[
g=\mathbb Pf,\qquad g_j=\partial_t^jg,\qquad
L=\nu\Delta,\qquad
B(v,w)=\mathbb P((v\cdot\nabla)w),
\]
\[
Q(v,w)=B(v,w)+B(w,v),\qquad
\mathcal K(u)=Lu-B(u,u),\qquad
\mathcal A_u h=Lh-Q(u,h).
\]
Here \(\mathcal K\) is the spatial differential expression in the
projected momentum equation, and \(\mathcal A_u\) is its derivative
with respect to the velocity argument. The Leray projection is
bounded in every Sobolev norm used below.

The momentum equation is \(V_1=\mathcal K(u)+g\). To differentiate it correctly, use
\[
\partial_t\mathcal K(u)=\mathcal A_uV_1,\qquad
\partial_t(\mathcal A_u h)=\mathcal A_u\partial_t h-Q(V_1,h).
\]
In particular, the coefficients of \(\mathcal A_u\) also change with time. Carrying this term gives
\[
\begin{aligned}
V_2&=\mathcal A_u\mathcal K(u)+S_2,\\
S_2&=\mathcal A_u g+g_1
=Lg-Q(u,g)+g_1,
\end{aligned}
\]
and then
\[
\begin{aligned}
V_3&=\mathcal A_u^2\mathcal K(u)
     -Q(\mathcal K(u),\mathcal K(u))+S_3,\\
S_3&=\mathcal A_u^2g+\mathcal A_u g_1
     -2Q(\mathcal K(u),g)-Q(g,g)+g_2.
\end{aligned}
\]
These are exact expansions of the full binomial rows. The coefficient two in the formula for \(S_3\) comes from differentiating both the velocity-dependent operator and the velocity argument. The quantities \(S_2,S_3\) collect source insertions into the velocity jet. They contain the prescribed source together with the velocity on which it acts.

The kinetic-energy inequality already bounds the first insertion
through the terminal time. Put
\[
U=\|u_0\|_2+\|f\|_{L^1_tL^2_x},\qquad
D^2=\nu^{-1}\left(\tfrac12\|u_0\|_2^2
                      +U\|f\|_{L^1_tL^2_x}\right).
\]
The actual forced history satisfies
\(\sup_{t<T}\|u(t)\|_2\le U\) and
\(\|\nabla u\|_{L^2_tL^2_x}\le D\).
Since \(B\) is followed by an \(L^2\)-contractive projection, the formula for \(S_2\) gives
\[
\begin{aligned}
\|S_2\|_{L^2_tL^2_x}
\le{}&\nu\|\Delta g\|_{L^2_tL^2_x}
      +\|g_1\|_{L^2_tL^2_x}\\
&+\|g\|_{L^\infty_{t,x}}D
 +\|\nabla g\|_{L^\infty_{t,x}}\sqrt{T}\,U
<\infty.
\end{aligned}
\]
Thus this first force--velocity insertion has a terminal bound independent of the critical velocity rectangle.

The next insertion admits an \(L^2_tH^{-3}_x\) estimate from
the kinetic bounds. The weaker spatial norm accommodates the
additional derivatives in this row. Define
\[
Z=\|u\|_{L^2_tH^1_x}\le(TU^2+D^2)^{1/2},\qquad
G=\max_{0\le j\le2}\sup_{t\le T}\|g_j(t)\|_{H^6}.
\]
The divergence form and Sobolev duality give
\[
\begin{aligned}
\|Q(u,g)\|_{L^2L^2}&\le C GZ,\\
\|B(u,u)\|_{L^2H^{-3/2}}&\le C UZ,\\
\|Q(u,Q(u,g))\|_{L^2H^{-3}}&\le C G UZ,\\
\|Q(Lu,g)\|_{L^2H^{-2}}&\le C\nu GZ,\\
\|Q(B(u,u),g)\|_{L^2H^{-5/2}}&\le C G UZ.
\end{aligned}
\]
For the second line use \(u\otimes u\in L^2_tL^{3/2}_x\), obtained from \(u\in L^\infty_tL^2_x\cap L^2_tL^6_x\), and the dual of \(H^{1/2}\hookrightarrow L^3\). For the third line, \(u\otimes Q(u,g)\) belongs to \(L^2_tL^1_x\); pairing its divergence against \(H^3\) is bounded because \(H^2\hookrightarrow L^\infty\). Smooth multiplication by \(g\) and its derivatives gives the final two lines. Expanding \(\mathcal A_u^2g\) in the formula for \(S_3\) now gives
\[
\|S_3\|_{L^2_tH^{-3}_x}
\le C\left[
\sqrt{T}\bigl((1+\nu+\nu^2)G+G^2\bigr)
+G(1+\nu+U)Z\right]<\infty.
\]
The constant \(C\) depends on fixed Sobolev product and multiplier constants. All velocity dependence is contained in the kinetic quantities \(U\) and \(Z\). The two estimates control \(S_2\) and \(S_3\) at the stated spatial heights.

The all-order version is also explicit. Define spatial jet polynomials by
\[
\mathcal J_0(v)=v,\qquad
\mathcal J_{n+1}(v)=D\mathcal J_n(v)[\mathcal K(v)].
\]
Here \(D\mathcal J_n(v)[h]\) means the derivative of the spatial polynomial with respect to \(v\) in direction \(h\). Along the actual forced history,
\[
V_n=\mathcal J_n(u)+S_n,\qquad S_0=0,
\]
\[
S_{n+1}=D\mathcal J_n(u)[g]+\partial_t S_n.
\]
This follows by the chain rule and \(u_t=\mathcal K(u)+g\). Every \(S_n\) is a finite differential expression with at least one source factor. The recurrence retains all velocity factors as well. The bounds above apply at the specified spatial heights.


\subsection{Time-integrated completion and terminal integrability}

To apply Proposition~\ref{r:prop:source-primitive} to the
completed-height identity \eqref{r:eq:completed}, take
\(r=0\), \(k\ge2\), \(m=k-2\), and \(H=E_{0,1/2}\). Define
\[
 \mathcal N_{s,k}:=\sum_{j=0}^{k-1}(-2\nu)^j
 \partial_t^{k-1-j}\mathcal P_{0,s+j}.
\]
Applying \(I^m\) to the full completion gives
\begin{equation}
 \begin{aligned}
 (-2\nu)^kI^m C^f_{0,k}
 &=H''-\sum_{\ell=0}^{m-1}H^{(\ell+2)}(0)\frac{t^\ell}{\ell!}\\
 &\quad-I^m\mathcal N_{1/2,k}-I^m\mathcal S_{1/2,k}.
 \end{aligned}
 \label{r:eq:integrated-completion}
\end{equation}
The source term on the last line has the proved finite bound. The energy
derivative and complete nonlinear term remain in the expression; a bound
on their combination has not been obtained from that source bound.

A full-interval Sobolev coordinate also specifies an integrability
requirement. For example, \(V_r\in L^2(0,T;\dot H^{k+1/2})\) means
\[
 \int_0^T C^f_{r,k}(t)\dd t
 =\frac12\int_0^T\norm{V_r(t)}_{\dot H^{k+1/2}}^2\dd t<\infty.
\]
This is the meaning of membership at that height. It is stronger than
the existence of every derivative on \([0,S]\) for each \(S<T\).
In particular, the full mixed product rule for \(W_{r,s}\) contains both
source and response factors, even when every source coordinate is finite.
For \(k\ge4\), boundedness of \(I^{k-2}C^f_{0,k}\) would itself give
only a weighted terminal integral,
\[
 \int_0^T(T-t)^{k-3}C^f_{0,k}(t)\dd t<\infty.
\]
The positive Volterra kernel gives this statement by monotone convergence;
its vanishing terminal weight does not give the unweighted integral above.
At \(k=2,3\), control of the full right side of
\eqref{r:eq:integrated-completion} would suffice. These distinctions specify which terminal integrals follow from
the integrated source estimate.
